(1) 位置算符 \(z=r\cos\theta\) 在空间反演下奇，因此同宇称态的对角矩阵元为零。角动量选择定则要求 \(\Delta l=\pm1,\Delta m=0\)，故 \(2s\) 仅与 \(2p,m=0\) 通过 \(z\) 一阶耦合；\(2p,m=\pm1\) 在这一四维空间中一阶矩阵元均零。此选择定则直接把四维问题简化成一个二维块和两个零块。

(2) 使用氢原子标准归一径向波函数 \(R_{20}=(2-r/a_0)e^{-r/(2a_0)}/(2\sqrt2a_0^{3/2})\)、\(R_{21}=(r/a_0)e^{-r/(2a_0)}/(2\sqrt6a_0^{3/2})\)，并配合 \(\int Y_{00}^*\cos\theta Y_{10}d\Omega=1/\sqrt3\)。积分 \(\int_0^\infty r^3R_{20}R_{21}dr\) 得 \(-3\sqrt3a_0\)，故 \(\langle2s|z|2p,0\rangle=-3a_0\)（相位约定可改变符号）。因此该简并子空间二维扰动块可写 \(e\mathcal E\begin{pmatrix}0&-3a_0\\-3a_0&0\end{pmatrix}\)。

(3) 对角化微扰二维块得到一阶能量修正 \(\Delta E_\pm=\pm3ea_0\mathcal E\)，而两个 \(m=\pm1\) 态仍为 \(\Delta E=0\)。对应归一化混合态为 \((|2s\rangle\mp|2p,0\rangle)/\sqrt2\)，具体正负号与前述矩阵选定相位有关。由于原能级严格轨道简并，微扰能量在这一理想模型中对场强线性，这和普通非简并态常见的二阶极化能量不同。

(4) 一阶不受影响的 \(m=\pm1\) 并不意味着它们在任意场强下没有二阶 Stark 位移，而只是所截取的简并子空间中一阶选择定则使其矩阵元为零。实际氢原子存在精细结构和 Lamb 位移，若这些能级分裂大于电场耦合 \(3ea_0\mathcal E\)，应先用真实非简并谱并作不同次序的微扰，能量对弱场可主要是二阶变化；不能无限外推理想简并线性公式。
